\section{Hai ràng buộc trên một ma trận trọng số}
\label{sec:ch10-hai-rang-buoc}

\index{thiên kiến quy nạp!hai ràng buộc của tích chập}%
Một tầng \tn{tích chập}{convolution} không phải một loại tầng mới. Nó là tầng
kết nối đầy đủ của chương~\ref{ch:ky-hieu-va-bptt} với hai ràng buộc đặt lên,
và hai ràng buộc ấy đặt được riêng rẽ. Bello và cộng sự viết ra gọn nhất, ở
ngay câu mở đầu bài của họ~\autocite{bello2019attention}:

\begin{quote}
\enquote{The design of the convolutional layer imposes 1) locality via a
limited receptive field and 2) translation equivariance via weight sharing.}
\end{quote}

Ràng buộc thứ nhất là \tn{tính cục bộ}{locality}: một unit đầu ra chỉ được nhìn
một lân cận nhỏ của đầu vào, gọi là \tn{trường tiếp nhận}{receptive field} của
nó, thay vì nhìn tất cả. Ràng buộc thứ hai là
\tn{chia sẻ trọng số}{weight sharing}: các unit đứng ở chỗ khác nhau nhưng nhìn
theo cùng một hướng thì phải dùng \emph{chung một} tham số.

Hình~\ref{fig:ch10-hai-rang-buoc} đặt ba chế độ cạnh nhau trên cùng một tầng
nhỏ, và chỗ đáng nhìn nhất ở đó là hai ngăn bên phải có \emph{cùng số cạnh}.
Chia sẻ trọng số không bỏ đi kết nối nào; nó chỉ bắt các kết nối dùng chung
tham số.

% The picture measures 367.99pt against this book's 441.02pt measure, so it has
% 73pt to spare. Measured with \savebox rather than judged from the built page,
% because a tikzpicture that overflows inside \centering raises no Overfull box
% and neither the log nor check-chapter.ps1 can see it (quyet dinh 67).
% Recheck this number before widening the picture by so much as a node.
%
% Note for whoever repeats the measurement: the savebox has to be named without
% digits in it. \newsavebox{\ch10boxa} parses as \ch followed by the letters
% 10boxa, and the \typeout then prints the width of some other box entirely -
% it read 5.475pt here, which looks like a plausible small number rather than
% like an error.
\begin{figure}[htbp]
  \centering
  % Convolution la fully-connected voi hai rang buoc ap chong len nhau: cuc bo
% (moi don vi ra chi noi voi vai lang gieng, khong phai toan bo dau vao) roi
% dung chung tham so (cac canh cung do lech trong cua so dung chung dung mot
% trong so). Ba tam trai sang phai - day dac, cuc bo, cuc bo + dung chung -
% dung chung mot day 7 don vi vao / 5 don vi ra (kernel 3, buoc nhay 1, mot
% kenh) de dem canh bang mat; anh 2D that se roi ngay thanh mang nhen.
% Diem chinh: tam B va C co DUNG MOT TAP CANH (ca hai 15 canh) va chi khac o
% so THAM SO. B coi moi canh la mot bien doc lap rieng (20 tham so, bias
% khong dung chung). C bat cac canh cung do lech trong cua so - dau/giua/cuoi
% cua so - dung chung dung mot trong so, nen chi con 4 tham so; vi vay canh
% cua C tach ba kieu net (lien dam / dut / cham nhat) lap lai giong het nhau
% o ca nam vi tri ra, con canh cua A/B dung mot kieu net xam dong nhat vi moi
% canh la mot tham so rieng, khong co gi de lap lai.
% Khong ve: padding, stride > 1, nhieu kenh, ham kich hoat - do la chi tiet
% cai dat, khong phai hai rang buoc dang noi cua hinh nay.
% Mono-safe: ba kieu net trong tam C phan biet bang net + do dam, khong mau.
% Do rong DA DO bang \savebox tai cho goi hinh: 367.99pt = 12.93cm, tren mot
% measure 441.02pt. Do la con so co gia tri; dung uoc luong lai bang mat.
\begin{tikzpicture}[
  >= Stealth,
  vao/.style        = {draw, circle, fill = black!8, minimum size = 3.6mm,
                        inner sep = 0pt},
  ra/.style         = {draw, circle, fill = white, minimum size = 3.6mm,
                        inner sep = 0pt},
  canhthuong/.style = {black!38, line width = 0.35pt},
  canhdau/.style    = {black!75, line width = 0.7pt},
  canhgiua/.style   = {black!75, line width = 0.7pt, densely dashed},
  canhcuoi/.style   = {black!45, line width = 1.0pt, densely dotted},
  khung/.style      = {black!45, densely dotted, line width = 0.5pt},
  nho/.style        = {font = \scriptsize},
  ten/.style        = {font = \bfseries\scriptsize},
  ghichu/.style     = {font = \footnotesize, black!60},
]

  \pgfmathsetmacro\OA{0}
  \pgfmathsetmacro\OB{4.4}
  \pgfmathsetmacro\OC{8.8}

  % ================= TẤM A: dày đặc (x0 = 0) =================
  \foreach \k in {1, 2, 3, 4, 5, 6, 7} {
    \node[vao] (a-in\k) at ({\OA+(\k-1)*0.5}, 0) {};
  }
  \foreach \j in {1, 2, 3, 4, 5} {
    \node[ra] (a-out\j) at ({\OA+0.5+(\j-1)*0.5}, 2.0) {};
  }
  \foreach \j in {1, 2, 3, 4, 5} {
    \foreach \k in {1, 2, 3, 4, 5, 6, 7} {
      \draw[canhthuong] (a-in\k) -- (a-out\j);
    }
  }

  \begin{scope}[on background layer]
    \draw[khung, rounded corners = 2pt]
      ({\OA-0.4}, -1.6) rectangle ({\OA+3.4}, 2.35);
  \end{scope}
  \node[ten] at ({\OA+1.5}, -0.55) {Dày đặc};
  \node[ghichu, align = center] at ({\OA+1.5}, -1.15)
    {(mọi cặp đều nối)\\35 cạnh, 40 tham số};

  % ================= TẤM B: cục bộ, không dùng chung (x0 = 4.4) =================
  \foreach \k in {1, 2, 3, 4, 5, 6, 7} {
    \node[vao] (b-in\k) at ({\OB+(\k-1)*0.5}, 0) {};
  }
  \foreach \j in {1, 2, 3, 4, 5} {
    \node[ra] (b-out\j) at ({\OB+0.5+(\j-1)*0.5}, 2.0) {};
  }
  \foreach \j in {1, 2, 3, 4, 5} {
    \pgfmathtruncatemacro{\ka}{\j}
    \pgfmathtruncatemacro{\kb}{\j+1}
    \pgfmathtruncatemacro{\kc}{\j+2}
    \draw[canhthuong] (b-in\ka) -- (b-out\j);
    \draw[canhthuong] (b-in\kb) -- (b-out\j);
    \draw[canhthuong] (b-in\kc) -- (b-out\j);
  }

  \begin{scope}[on background layer]
    \draw[khung, rounded corners = 2pt]
      ({\OB-0.4}, -1.6) rectangle ({\OB+3.4}, 2.35);
  \end{scope}
  \node[ten] at ({\OB+1.5}, -0.55) {Cục bộ};
  \node[ghichu, align = center] at ({\OB+1.5}, -1.15)
    {(mỗi cạnh 1 tham số riêng)\\15 cạnh, 20 tham số};

  % ================= TẤM C: cục bộ + dùng chung, tích chập (x0 = 8.8) =================
  \foreach \k in {1, 2, 3, 4, 5, 6, 7} {
    \node[vao] (c-in\k) at ({\OC+(\k-1)*0.5}, 0) {};
  }
  \foreach \j in {1, 2, 3, 4, 5} {
    \node[ra] (c-out\j) at ({\OC+0.5+(\j-1)*0.5}, 2.0) {};
  }
  % Cùng đúng 15 cạnh của tấm B, nhưng mỗi cạnh mang kiểu nét theo ĐỘ LỆCH
  % của nó trong cửa sổ (đầu/giữa/cuối) - và kiểu đó lặp lại giống hệt nhau
  % ở cả năm vị trí ra: đó chính là dùng chung tham số.
  \foreach \j in {1, 2, 3, 4, 5} {
    \pgfmathtruncatemacro{\ka}{\j}
    \pgfmathtruncatemacro{\kb}{\j+1}
    \pgfmathtruncatemacro{\kc}{\j+2}
    \draw[canhdau]  (c-in\ka) -- (c-out\j);
    \draw[canhgiua] (c-in\kb) -- (c-out\j);
    \draw[canhcuoi] (c-in\kc) -- (c-out\j);
  }

  \begin{scope}[on background layer]
    \draw[khung, rounded corners = 2pt]
      ({\OC-0.4}, -1.6) rectangle ({\OC+3.4}, 2.35);
  \end{scope}
  \node[ten] at ({\OC+1.5}, -0.55) {Cục bộ + dùng chung};
  \node[ghichu, align = center] at ({\OC+1.5}, -1.15)
    {(mỗi độ lệch 1 trọng số)\\15 cạnh, 4 tham số};

  % ================= chú giải chung, đặt dưới cả ba tấm =================
  \node[nho, anchor = north, align = left, black!70, text width = 12cm]
    at ({\OB+1.9}, -1.85)
    {Ba tấm dùng chung một cấu trúc: 7 đơn vị vào (hàng dưới), 5 đơn vị ra
     (hàng trên), cửa sổ rộng 3, bước nhảy 1, một kênh. Dày đặc nối mọi cặp:
     35 cạnh, 35 trọng số + 5 bias = 40 tham số. Cục bộ giữ đúng 15 cạnh
     đó nhưng coi mỗi cạnh là một tham số riêng: 15 trọng số + 5 bias =
     20 tham số. Cục bộ + dùng chung giữ y nguyên 15 cạnh ấy, nhưng các cạnh
     cùng độ lệch trong cửa sổ - nét liền đậm là đầu cửa sổ, nét đứt là
     giữa, nét chấm nhạt là cuối - dùng lại đúng một trọng số ở mọi vị trí
     ra: chỉ còn 3 trọng số + 1 bias = 4 tham số.};

\end{tikzpicture}

  \caption{Hai ràng buộc, đặt lần lượt, trên một tầng 7 vào và 5 ra với lân cận
  cỡ 3. Ngăn giữa và ngăn phải có đúng cùng một tập cạnh: khác biệt duy nhất là
  ở ngăn phải, các cạnh cùng một độ lệch bị buộc phải dùng chung một tham số,
  và điều đó đưa 20 tham số xuống còn 4. Số cạnh không đổi, nên chi phí tính
  toán không đổi.}
  \label{fig:ch10-hai-rang-buoc}
\end{figure}

\subsection*{Mỗi ràng buộc đáng bao nhiêu}

\index{thiên kiến quy nạp!giá của tính cục bộ và chia sẻ}%
Đếm được, nên đếm. Bảng dưới lấy đúng tầng đầu tiên của mạng
mục~\ref{sec:ch10-anh-that} huấn luyện: ảnh 32x32 ba kênh vào, 32 kênh ra, lân
cận 3x3.

\begin{measured}{cùng một tầng dưới ba chế độ}
\begin{minted}{text}
regime  connections  weights      biases  parameters   conn/param
dense   100,696,064  100,663,296  32,768  100,696,064        1.00
local   917,504      884,736      32,768  917,504            1.00
conv    917,504      864          32      896             1024.00
\end{minted}

\begin{minted}{text}
each constraint as a divisor on the parameter count
  locality alone          109.75
  sharing, on top        1024.00
  both together        112384.00
\end{minted}

Đo bằng \code{experiments/ch10\_counts.py} tại tag \code{ch10}.
\end{measured}

Hai điều đọc ra.

Thứ nhất, hai ràng buộc \emph{không} ngang nhau. Tính cục bộ chia số tham số
cho 109.75, chia sẻ trọng số chia tiếp cho 1024. Cái thứ hai lớn hơn cái thứ
nhất gần mười lần.

Thứ hai, và đây mới là chỗ đáng nhớ: \textbf{con số 1024 đúng bằng số vị trí
không gian của tầng ấy, 32 nhân 32.} Không phải xấp xỉ. Chia sẻ trọng số chia
số tham số cho đúng số chỗ mà tầng ấy nhìn, vì đó chính là định nghĩa của nó -
mọi vị trí dùng lại một bộ trọng số. Tính cục bộ thì chia cho tỷ số hai trường
tiếp nhận, và tỷ số ấy phụ thuộc kích thước ảnh theo cách khác hẳn.

Hệ quả thực hành xuất hiện lại ở mục~\ref{sec:ch10-anh-that}, và nó khó chịu
hơn vẻ ngoài: cùng một ngân sách tham số mua cho một mạng tích chập nhiều hơn
hẳn so với mua cho một mạng không chia sẻ. Nên một phép so \enquote{cùng số
tham số} giữa hai kiến trúc ấy không phải một phép so công bằng theo nghĩa
người ta hay tưởng.

\subsection*{Dựng lại bài 1989}

\index{thiên kiến quy nạp!dựng lại số tham số LeCun 1989}%
Quyết định của sách này là một con số lấy từ bảng của bài phải được dựng lại từ
các bảng khác của chính bài trước khi in ra. Chương~\ref{ch:tien-huan-luyen}
chạy phép ấy trên bốn bài và ba bài không quay về. Ở đây tôi chạy lại vì một
quy tắc chỉ chạy khi đoán trước là sẽ hỏng thì không phải quy tắc.

Bài của LeCun và cộng sự năm 1989 là chỗ chia sẻ trọng số đi vào một mạng
nơ-ron thị giác được huấn luyện bằng
\tn{lan truyền ngược}{backpropagation}~\autocite{lecun1989zipcode}. Bài không
nhận là đã nghĩ ra nó - mục 3.2 dẫn về Rumelhart và cộng sự 1986 - nhưng nó là
chỗ định nghĩa gọn nhất, trong đúng một câu:

\begin{quote}
\enquote{Weight sharing was described in Rumelhart et al. (1986) for the
so-called T-C problem and consists in having several connections (links)
controlled by a single parameter (weight). It can be interpreted as imposing
equality constraints among the connection strengths.}
\end{quote}

\enquote{Equality constraints among the connection strengths} là ràng buộc thứ
hai, viết ra bằng lời, năm 1989.

\begin{measured}{mạng của LeCun 1989, dựng lại từ chính mục 3.3 của bài}
\begin{minted}{text}
layer   connections  weights  biases  parameters  conn/param
H1      19,968       300      768     1,068          18.697
H2      38,592       2,400    192     2,592          14.889
H3      5,790        5,760    30      5,790           1.000
output  310          300      10      310             1.000
TOTAL   64,660       8,760    1,000   9,760           6.625
\end{minted}

Bài in: 1256 unit, 64,660 kết nối, 9,760 tham số độc lập.

Đo bằng \code{experiments/ch10\_counts.py} tại tag \code{ch10}.
\end{measured}

Khớp cả hai tổng, đúng tới đơn vị. Tỷ số là 6.625 chẵn.

Chuyện đáng đọc nằm ở hàng H1, và nó là chuyện bài nói ra nhưng không nhấn.
Cùng mục 3.3:

\begin{quote}
\enquote{Units do not share their biases (thresholds). Each unit thus has 25
input lines plus a bias.}
\end{quote}

Nên trong 1,068 tham số tự do của tầng tích chập đầu tiên trong lịch sử, có
\emph{768 là bias} và chỉ 300 là trọng số kernel.

\begin{measured}{H1 tách ra làm hai phần}
\begin{minted}{text}
  H1 shared weights                300
  H1 unshared biases               768
  biases as a share of H1       0.7191
\end{minted}

Đo bằng \code{experiments/ch10\_counts.py} tại tag \code{ch10}.
\end{measured}

71.91\%. Chia sẻ trọng số bỏ đi các trọng số và không đụng tới bias, nên cái
còn lại chủ yếu là bias. Chuyện này không sai ở đâu cả, và nó không đổi luận
điểm nào của bài; tôi in ra vì nó là thứ mà một người đọc chữ \enquote{chỉ 1068
tham số} sẽ không đoán được, và vì nó cho thấy phép chia 1024 ở bảng trước chỉ
áp lên một nửa của một tầng.

\subsection*{Dựng lại LeNet-5, và chỗ nó không khớp}

\index{thiên kiến quy nạp!dựng lại LeNet-5}%
Chín năm sau, cùng nhóm tác giả, bài đặt tên cho ba ý tưởng cùng một
lúc~\autocite{lecun1998gradient}, mục II.A:

\begin{quote}
\enquote{Convolutional Networks combine three architectural ideas to ensure
some degree of shift, scale, and distortion invariance: local receptive fields,
shared weights (or weight replication), and spatial or temporal sub-sampling.}
\end{quote}

Ý tưởng thứ ba, \tn{giảm mẫu}{subsampling}, là thứ mục~\ref{sec:ch10-dang-bien}
sẽ cho thấy phải trả giá.

\begin{measured}{LeNet-5, dựng lại từ mục II.B}
\begin{minted}{text}
layer   connections  parameters
C1      122,304      156
S2      5,880        12
C3      151,600      1,516
S4      2,000        32
C5      48,120       48,120
F6      10,164       10,164
output  840          0
TOTAL   340,908      60,000
\end{minted}

Bài in: 340,908 kết nối, \enquote{only 60,000 trainable free parameters}.

Đo bằng \code{experiments/ch10\_counts.py} tại tag \code{ch10}.
\end{measured}

\textbf{60,000 là con số tròn nhưng chính xác.} Nó là tổng, không phải một phép
làm tròn, và tôi đã chờ đợi điều ngược lại.

Hàng \code{output} là hàng phải đi tìm. Cộng sáu tầng đầu thì tham số vẫn đúng
60,000 còn kết nối chỉ ra 340,068, thiếu đúng 840. LeNet-5 kết thúc bằng 10
unit hàm cơ sở bán kính - radial basis function, viết tắt là RBF - mỗi unit 84
đầu vào; trọng số của chúng bài đặt bằng tay và không huấn luyện. Nên tầng ấy đóng góp 840 kết nối và \emph{không} tham số tự do nào,
và nó vào một tổng mà không vào tổng kia. Chỗ lệch 840 là cách tôi tìm ra nó.

\subsection*{Chia sẻ trọng số không làm mạng nhỏ đi}

\index{thiên kiến quy nạp!phần lớn tham số nằm ở tầng không chia sẻ}%
Còn một cột nữa trong hai bảng trên, và nó nói một chuyện ngược với câu người
ta hay kể.

\begin{measured}{tham số nằm ở đâu, trong hai mạng tích chập kinh điển}
\begin{minted}{text}
LeCun 1989
  H1 + H2, the shared layers     3,660
  H3 + output, fully connected   6,100
  fully connected share         0.6250

LeNet-5
  C1..S4, the shared layers      1,716
  C5 + F6, fully connected      58,284
  fully connected share         0.9714
\end{minted}

Đo bằng \code{experiments/ch10\_counts.py} tại tag \code{ch10}.
\end{measured}

62.50\% ở bài 1989, và \textbf{97.14\% ở LeNet-5}. Kiến trúc đặt tên cho chia
sẻ trọng số tiêu 97\% tham số của nó ở những tầng chia sẻ đúng con số không.

Cách đọc đúng không phải là bài mâu thuẫn với chính nó. Là chia sẻ trọng số
không làm \emph{mạng} nhỏ đi; nó làm \emph{bộ trích đặc trưng} nhỏ đi. Phần
tích chập của LeNet-5 nặng 1,716 tham số và sinh ra một biểu diễn mà một bộ
phân loại kết nối đầy đủ nhỏ đọc được. Đó là phân công lao động, không phải một
phép nén.

Và nó là lý do câu hỏi của chương này không trả lời được bằng cách đếm tham số.
Cái mà hai ràng buộc mua được không phải kích thước. Mục sau là chỗ nói nó mua
được cái gì.
